%% 4-14 外積：回す力と右ねじ（第13回：組 V4-7。ページ制限なし。偶数ページ）
%% 物理シリーズ（力学）のために必ず書く形：力のモーメント（トルク）、平面の M＝xF_y−yF_x、
%%   外積の向き（右ねじ）と大きさ（平行四辺形の面積）、トルク r×F、角運動量、q v×B、フレミングの左手の法則
\makeatletter
% 斜めから見た空間の図（4-13 と同じかき方）
\tikzset{ivo3d/.style={x={({-0.49*#1},{-0.34*#1})}, y={(#1,0)}, z={(0,#1)}}}
% 式を、まん中に1行で置く
\newcommand\bkivoline[1]{\par\addvspace{3pt}{\centering #1\par}\addvspace{3pt}\noindent\ignorespaces}
% 説明の段の見出し（丸数字＋太字の1行）
\newcommand\bkivostep[2]{\par\addvspace{6pt}\Needspace{5\baselineskip}\noindent
  {\bkcircled[bkOrange]{#1}\enspace\sffamily\bfseries\color{bkOrange!75!black}#2}\par\nobreak\vspace{2pt}}
% 小見出し（それだけで1行）
\newcommand\bkivohead[1]{\par\addvspace{4pt}\Needspace{6\baselineskip}\noindent\textbf{#1}\par\nobreak}
% 3本の軸
\newcommand\bkivoaxes[3]{%
  \draw[bkGray, line width=0.6pt, -{Stealth[length=4pt]}] (0,0,0) -- (#1,0,0) node[anchor=east, inner sep=2pt, font=\sffamily\small, text=bkInk] {$x$};
  \draw[bkGray, line width=0.6pt, -{Stealth[length=4pt]}] (0,0,0) -- (0,#2,0) node[anchor=west, inner sep=2pt, font=\sffamily\small, text=bkInk] {$y$};
  \draw[bkGray, line width=0.6pt, -{Stealth[length=4pt]}] (0,0,0) -- (0,0,#3) node[anchor=south, inner sep=2pt, font=\sffamily\small, text=bkInk] {$z$};}
% 紙面から手前へ出る向き ⊙ と、紙面のおくへ入る向き ⊗（canvas の点 #1、半径 #2）
\newcommand\bkivoout[2]{\draw[bkInk, line width=0.8pt, fill=white] #1 circle (#2); \fill[bkInk] #1 circle (#2*0.28);}
\newcommand\bkivoin[2]{\draw[bkInk, line width=0.8pt, fill=white] #1 circle (#2);
  \draw[bkInk, line width=0.8pt] #1 ++(45:#2) -- ++(225:2*#2) #1 ++(135:#2) -- ++(-45:2*#2);}
% 指（4-13 と同じ形）
\newcommand\bkivofinger[3]{%
  \draw[bkOrange!70!black, line width=#3, line cap=round] #1 -- #2;
  \draw[bkOrange!22, line width=#3-1pt, line cap=round] #1 -- #2;}
% N·m（中の点つき）
\newcommand\bkivoNm[1]{\qty[inter-unit-product=\ensuremath{\cdot}]{#1}{N.m}}
\makeatother

\bkchapter{4}{14}{外積：回す力と右ねじ}{\bklevelrange{4}{5}}{外積：回す力と右ねじ}
\bkmichishirube{4-7,4-12,4-13}{7-4,8-10}

\begin{goal}
\begin{itemize}[label={\color{\bkTT}●}, leftmargin=1.6\zw, itemsep=0pt, topsep=1pt]
\item \bklevel{3}\enspace 回す働き（力のモーメント）を「力×うでの長さ」で求め、平面では $M=xF_y-yF_x$ で計算できる。
\item \bklevel{4}\enspace 外積の向き（右ねじ）と大きさ（平行四辺形の面積）がわかり、成分で計算できる。
\item \bklevel{5}\enspace トルク $\vec r\times\vec F$ や、磁場から受ける力の向きを、外積で求められる。
\end{itemize}
\end{goal}

\begin{hatena}
\noindent\begin{minipage}[c]{0.56\linewidth}
ドアのノブ（取っ手）は、ちょうつがいから遠い所についています。ちょうつがいのすぐそばを押して、ドアを開けようとしてみましょう。とても重く感じます。

同じドアなのに、押す場所で重さがちがうのはなぜかな？ ノブがちょうつがいのそばにあったら、どうなるかな？
\end{minipage}\hfill
\begin{minipage}[c]{0.4\linewidth}
\centering
\begin{tikzpicture}[x=1mm, y=1mm]
  % 真上から見たドア
  \fill[bkGray!40] (-4,-3) rectangle (0,3);
  \fill[bkOrange!30] (0,-1.5) rectangle (40,1.5);
  \draw[bkOrange!70!black, line width=0.6pt] (0,-1.5) rectangle (40,1.5);
  \fill[bkInk] (0,0) circle (1.3);
  \fill[bkYellow!80!black] (36,1.5) rectangle (38,4);
  \draw[bk arrow, bkBlue, line width=1.3pt] (36,13) -- (36,3);
  \draw[bk arrow, bkRed, line width=1.3pt] (5,13) -- (5,3);
  \node[font=\sffamily\footnotesize, anchor=north] at (0,-4) {ちょうつがい};
  \node[font=\sffamily\footnotesize, anchor=north] at (37,-4) {ノブ};
  \node[font=\sffamily\footnotesize, text=bkRed!80!black, anchor=south] at (5,13.5) {重い};
  \node[font=\sffamily\footnotesize, text=bkBlue!80!black, anchor=south] at (36,13.5) {軽い};
  \node[font=\sffamily\footnotesize, text=bkGray, anchor=north] at (20,-8.5) {（真上から見た図）};
\end{tikzpicture}
\end{minipage}
\end{hatena}

\noindent ものを「動かす」働きは、力の大きさと向きで決まりました（4-10・4-11）。この章では、ものを「回す」働きを数で表します。回す働きには、力の大きさのほかに、「どこを押すか」が大切です。そして、回す働きを向きのある量として表すと、ベクトルの新しいかけ算「外積」が出てきます。

%% ---------------- やってみよう・思い出そう ----------------
\begin{yattemiyou}[ドアを開ける力をはかろう]
\noindent\begin{minipage}[c]{0.58\linewidth}
\textbf{用意するもの}\par
ばねばかり（またはキッチンスケール）、メジャー、ひも
\begin{enumerate}[label={\bkcircled[bkGreen]{\arabic*}}, leftmargin=2.4\zw, itemsep=1pt, topsep=3pt]
\item ドアに、ちょうつがいから\qty{0.1}{m}、\qty{0.2}{m}、\qty{0.4}{m}、\qty{0.6}{m}、\qty{0.75}{m}の所に、ふせんで印をつけよう。
\item 印の所にひもをはり、ドアの面に直角の向きに、ばねばかりでゆっくり引こう。
\item ドアが動き始めたときの力を読もう。
\end{enumerate}
\end{minipage}\hfill
\begin{minipage}[c]{0.38\linewidth}
\centering
\begin{tikzpicture}[x=1mm, y=1mm]
  \fill[bkGray!40] (-3,-3) rectangle (0,3);
  \fill[bkOrange!30] (0,-1.5) rectangle (42,1.5);
  \draw[bkOrange!70!black, line width=0.6pt] (0,-1.5) rectangle (42,1.5);
  \fill[bkInk] (0,0) circle (1.2);
  \foreach \p in {5.6,11.2,22.4,33.6} {\fill[bkYellow!80!black] (\p-0.8,1.5) rectangle (\p+0.8,2.6);}
  \draw[bkInk, line width=0.5pt] (33.6,1.5) -- (33.6,11);
  \filldraw[fill=white, draw=bkInk, line width=0.6pt, rounded corners=1mm] (31.6,11) rectangle (35.6,21);
  \draw[bkInk, line width=0.4pt] (32.6,13) -- (34.6,13) (32.6,15) -- (34.6,15) (32.6,17) -- (34.6,17);
  \node[font=\sffamily\footnotesize, anchor=east] at (30.6,16) {ばねばかり};
  \draw[bkGray, line width=0.4pt, <->] (0,-5) -- (33.6,-5);
  \node[font=\sffamily\footnotesize, anchor=north, fill=bkGreenBg, inner sep=1pt] at (16.8,-5.6) {ちょうつがいからの長さ};
\end{tikzpicture}
\end{minipage}
\par\vspace{6pt}\noindent ミナさんの家のドアではかった結果です。
\begin{bkcenter}
\begin{bktable}{colspec={Q[l,m,wd=40mm] *{5}{Q[c,m,wd=11mm]}}, rowsep=1.5pt}
ちょうつがいからの長さ（m） & 0.1 & 0.2 & 0.4 & 0.6 & 0.75 \\
動き始めたときの力（N） & 6.0 & 3.0 & 1.5 & 1.0 & 0.8 \\
力×長さ & & & & & \\
\end{bktable}
\end{bkcenter}
\end{yattemiyou}

\noindent 表の3行めを計算してみましょう。$6.0\times0.1=0.6$、$3.0\times0.2=0.6$、$1.5\times0.4=0.6$、$1.0\times0.6=0.6$、$0.8\times0.75=0.6$。押す場所を変えると力は変わりますが、\textbf{力×長さ}は、どこでも同じ0.6になりました。長さが2倍になると、力は $\dfrac{1}{2}$ 倍です。2-15で学んだ反比例の形です。

\begin{omoidasou}
\bkrecall{2-15}{てこは、（おもりの重さ）×（支点からの距離）が左右で等しいときにつり合う。}
\bkrecall{4-7}{三角形の面積は $\dfrac{1}{2}ab\sin C$（2辺とそのあいだの角から）。}
\bkrecall{4-11}{力 $F$ を、$x$ 軸と角 $\theta$ の向きに加えると、成分は $(F\cos\theta,\ F\sin\theta)$。}
\bkrecall{4-12}{内積 $\vec a\cdot\vec b=|\vec a||\vec b|\cos\theta$。答えはふつうの数（スカラー）。}
\bkrecall{4-13}{空間の点は $(x,\ y,\ z)$。軸は右手系に並べる（$x$ から $y$ へ回すと、右ねじは $z$ の向きに進む）。}
\end{omoidasou}

\noindent この章では、$\approx$（「ほぼ等しい」というしるし。5-8でくわしく）も使います。

%% ---------------- 数学の道具 ----------------
\bksection{doguu}{数学の道具}

\bksubsection{1}{回す働き＝力×うでの長さ}

\leavevmode\label{bk:4-14:moment}やってみようでは、どこを押しても「力×長さ」が0.6でした。ドアが動き始めるのは、この「力×長さ」がある大きさ（ここでは0.6）になったときだ、と考えられます。

\noindent\begin{minipage}[c]{0.5\linewidth}
\centering
\begin{tikzpicture}[x=1mm, y=1mm]
  \fill[bkGray!40] (-4,-3) rectangle (0,3);
  \fill[bkOrange!30] (0,-1.5) rectangle (56,1.5);
  \draw[bkOrange!70!black, line width=0.6pt] (0,-1.5) rectangle (56,1.5);
  \fill[bkInk] (0,0) circle (1.3);
  \draw[bk arrow, bkRed, line width=1.6pt] (8,24) -- (8,3);
  \draw[bk arrow, bkBlue, line width=1.2pt] (52,8) -- (52,3);
  \draw[bkGray, line width=0.4pt, <->] (0,-6) -- (8,-6);
  \draw[bkGray, line width=0.4pt, <->] (0,-15) -- (52,-15);
  \node[font=\sffamily\footnotesize, anchor=north, inner sep=1pt] at (4,-6.8) {0.1 m};
  \node[font=\sffamily\footnotesize, anchor=north, inner sep=1pt] at (26,-15.8) {0.75 m};
  \node[font=\sffamily\footnotesize, text=bkRed!80!black, anchor=west] at (9.5,18) {6.0 N};
  \node[font=\sffamily\footnotesize, text=bkBlue!80!black, anchor=east] at (50.5,8) {0.8 N};
  \draw[bkInk, line width=0.8pt, -{Stealth[length=4pt]}] (-6:22) arc (-6:-38:22);
  \node[font=\sffamily\footnotesize, anchor=west, inner sep=1.5pt] at (-24:22.5) {ドアが回る};
\end{tikzpicture}
\end{minipage}\hfill
\begin{minipage}[c]{0.46\linewidth}
ちょうつがいのように、回る中心になる所を\textbf{回転の中心}（支点）といい、中心から力をかける点までの長さを\textbf{うでの長さ}といいます。

回す働きは、
\bkivoline{\fbox{力×うでの長さ}}
で表せます。
\end{minipage}

\par\vspace{4pt}
この「力×うでの長さ」を、\textbf{力のモーメント}といいます。\bkidx{ちからのもーめんと}{力のモーメント}機械の分野では\textbf{トルク}ともいいます。力の単位はN（ニュートン）、長さの単位はmなので、力のモーメントの単位は\textbf{N·m}（ニュートンメートル）です。やってみようのドアは、$\qty{6.0}{N}\times\qty{0.1}{m}=\bkivoNm{0.6}$ の力のモーメントで動き始めました。

\bkivohead{\bkhatenaref の答え}
ドアを回すには、決まった大きさ（この家のドアでは\bkivoNm{0.6}）の力のモーメントが必要です。力のモーメントは「力×うでの長さ」なので、うでが長い（ちょうつがいから遠い）ほど、小さい力ですみます。ノブを\qty{0.75}{m}の所につけると\qty{0.8}{N}で開きますが、\qty{0.1}{m}の所だと\qty{6.0}{N}、7.5倍の力がいります。だから、ノブはちょうつがいから遠い所についているのです。

\begin{seikaku}
\textbf{力がはたらく点（作用点）が大切}\par
4-10では、「力がどこにはたらくか（作用点）と、物を回す働きは、4-14で調べる」と予告しました。物を動かす（ずらす）働きだけなら、力の大きさと向きで決まりました。でも、回す働きは、同じ大きさ・同じ向きの力でも、作用点が回転の中心からどれだけはなれているかで変わります。ドアのノブを押すのと、ちょうつがいのそばを押すのとでは、力のベクトルは同じでも、回す働きは7.5倍ちがいました。だから、回す働きを考えるときは、力のベクトルに加えて、作用点の位置（回転の中心から作用点へのベクトル）も必要になります。この2つから回す働きを作る計算が、⑤の外積です。
\end{seikaku}

\bkivostep{1}{つり合い：左回りと右回りが等しい}
シーソーで考えます。左の人は、シーソーを左回り（反時計回り）に回そうとします。右の人は、右回り（時計回り）に回そうとします。2つの力のモーメントが等しいとき、シーソーはどちらにも回らず、つり合います。これは2-15のてこのつり合いと同じです。

\begin{bkfigure}
\begin{tikzpicture}[x=1mm, y=1mm]
  \fill[bkGray!50] (0,0) -- (-5,-9) -- (5,-9) -- cycle;
  \fill[bkOrange!35] (-46,0) rectangle (46,2.2);
  \draw[bkOrange!70!black, line width=0.5pt] (-46,0) rectangle (46,2.2);
  \bkperson{(-36,2.2)}{16}{bkBlue}
  \bkperson{(27,2.2)}{18}{bkRed}
  \draw[bk arrow, bkBlue, line width=1.4pt] (-36,-1) -- (-36,-12);
  \draw[bk arrow, bkRed, line width=1.4pt] (27,-1) -- (27,-14);
  \node[font=\sffamily\footnotesize, text=bkBlue!80!black, anchor=north] at (-36,-12.5) {\qty{300}{N}};
  \node[font=\sffamily\footnotesize, text=bkRed!80!black, anchor=north] at (27,-14.5) {\qty{400}{N}};
  \draw[bkGray, line width=0.4pt, <->] (-36,-22) -- (0,-22);
  \draw[bkGray, line width=0.4pt, <->] (0,-22) -- (27,-22);
  \node[font=\sffamily\footnotesize, anchor=north, inner sep=1.5pt] at (-18,-22.4) {\qty{1.2}{m}};
  \node[font=\sffamily\footnotesize, anchor=north, inner sep=1.5pt] at (13.5,-22.4) {\qty{0.9}{m}};
  \draw[bkGray, line width=0.3pt, dashed] (0,-9) -- (0,-22);
  \node[font=\sffamily\footnotesize, align=left, anchor=west] at (50,4) {左回り：$300\times1.2=360$\\右回り：$400\times0.9=360$};
\end{tikzpicture}
\bkfigcaption{重力（重さ）\qty{300}{N}の人が支点から\qty{1.2}{m}、\qty{400}{N}の人が\qty{0.9}{m}にすわると、左回りと右回りの力のモーメントが等しく、つり合う。}
\end{bkfigure}

\bksubsection{2}{ななめに押すと：うでに直角な部分だけが回す}

\leavevmode\label{bk:4-14:naname}ドアを押すとき、ドアの面に直角ではなく、ななめに押したらどうなるでしょう。極端な場合を考えます。ドアのノブを、ちょうつがいに向かってまっすぐ押すと、ドアはまったく回りません。力はちょうつがいを押しつけるだけです。

ななめの力は、2通りの見方で考えられます。

\begin{bkfigure}
\begin{tikzpicture}[x=0.92mm, y=0.92mm]
  % 左：力を2つに分ける
  \fill[bkInk] (0,0) circle (1.2);
  \draw[bkOrange!70!black, line width=2.4pt] (0,0) -- (40,0);
  \draw[bk arrow, bkRed, line width=1.6pt] (40,0) -- ({40+20*cos(30)},{20*sin(30)});
  \draw[bk arrow, bkBlue, line width=1.2pt] (40,0) -- ({40+20*cos(30)},0);
  \draw[bk arrow, bkGreen!70!black, line width=1.2pt] (40,0) -- (40,{20*sin(30)});
  \draw[bkGray, line width=0.4pt, dashed] (40,{20*sin(30)}) -- ({40+20*cos(30)},{20*sin(30)}) -- ({40+20*cos(30)},0);
  \draw[bkInk, line width=0.5pt] (47,0) arc (0:30:7);
  \node[font=\sffamily\footnotesize] at (49.8,2.4) {$\theta$};
  \node[font=\sffamily\footnotesize, text=bkRed!80!black, anchor=south west, inner sep=1pt] at ({40+20*cos(30)},{20*sin(30)}) {$F$};
  \node[font=\sffamily\footnotesize, text=bkGreen!50!black, anchor=east, inner sep=1.5pt] at (40,6) {$F\sin\theta$};
  \node[font=\sffamily\footnotesize, text=bkBlue!80!black, anchor=north, inner sep=1.5pt] at (49,-1) {$F\cos\theta$};
  \node[font=\sffamily\footnotesize, anchor=north, inner sep=1.5pt] at (20,-1.5) {うで $r$};
  \node[font=\sffamily\footnotesize, anchor=north] at (0,-2) {中心};
  \node[font=\sffamily\small\bfseries, anchor=south] at (32,16) {見方1：力を分ける};
  % 右：うでの長さを作用線までの距離にする
  \begin{scope}[shift={(78,0)}]
  \fill[bkInk] (0,0) circle (1.2);
  \draw[bkOrange!70!black, line width=2.4pt] (0,0) -- (40,0);
  \draw[bkGray, line width=0.5pt, dashed] ({40-37*cos(30)},{-37*sin(30)}) -- ({40+12*cos(30)},{12*sin(30)});
  \draw[bk arrow, bkRed, line width=1.6pt] (40,0) -- ({40+20*cos(30)},{20*sin(30)});
  \draw[bkPurple, line width=1.3pt] (0,0) -- ({40*sin(30)*sin(30)},{-40*sin(30)*cos(30)});
  \draw[bkInk, line width=0.4pt] ($({40*sin(30)*sin(30)},{-40*sin(30)*cos(30)})+({cos(30)*2.5},{sin(30)*2.5})$) -- ++({sin(30)*2.5},{-cos(30)*2.5}) -- ++({-cos(30)*2.5},{-sin(30)*2.5});
  \node[font=\sffamily\footnotesize, text=bkPurple, anchor=east, inner sep=1.5pt] at (4.5,-8) {$r\sin\theta$};
  \node[font=\sffamily\footnotesize, anchor=south, inner sep=1.5pt] at (20,0.8) {$r$};
  \node[font=\sffamily\footnotesize, text=bkRed!80!black, anchor=south west, inner sep=1pt] at ({40+20*cos(30)},{20*sin(30)}) {$F$};
  \node[font=\sffamily\small\bfseries, anchor=south] at (30,16) {見方2：うでをはかり直す};
  \end{scope}
\end{tikzpicture}
\bkfigcaption{見方1：うでにそった部分 $F\cos\theta$（青）は中心を引くだけで、回さない。うでに直角な部分 $F\sin\theta$（緑）だけが回す。見方2：力の矢印をのばした線（点線）に、中心から直角に引いた長さ $r\sin\theta$（むらさき）を、うでの長さと考える。}
\end{bkfigure}

\begin{itemize}[leftmargin=1.2\zw, itemsep=3pt, topsep=2pt, label={\color{\bkTT}●}]
\item \textbf{見方1：力を分ける}\quad 力 $F$ を、うでにそった部分と、うでに直角な部分に分けます（4-11の力の分解）。うでとのなす角を $\theta$ とすると、うでに直角な部分は $F\sin\theta$。回す働きは $r\times F\sin\theta$。
\item \textbf{見方2：うでをはかり直す}\quad 力の矢印をまっすぐのばした線（作用線）に、回転の中心から垂線を下ろします。その長さは $r\sin\theta$（直角三角形で、斜辺 $r$ の向かいの辺）。回す働きは $r\sin\theta\times F$。
\end{itemize}
どちらで考えても、
\bkivoline{\fbox{力のモーメント $M=rF\sin\theta$\qquad（$\theta$ は、うでと力のなす角）}}
になります。$\theta=90°$（うでに直角に押す）なら $\sin90°=1$ で、いちばん大きい $rF$。$\theta=0°$（中心に向かって押す）なら $\sin0°=0$ で、回しません。

\bksubsection{3}{平面で、成分から計算する：$M=xF_y-yF_x$}

\leavevmode\label{bk:4-14:heimen}力が座標で与えられたときは、角度を求めなくても計算できる方法があります。回転の中心を原点Oにして、力 $\vec F=(F_x,\ F_y)$ が点P$(x,\ y)$ にかかっているとします。

\bkivostep{1}{回る向きに、符号をつける}
回す働きには、左回りと右回りの2つの向きがあります。4-8の角と同じように、\textbf{左回り（反時計回り）を＋、右回り（時計回り）を－}と決めます。シーソーなら、左の人の力のモーメントは＋、右の人のは－です。つり合いは「合計が0」と言いかえられます。

\bkivostep{2}{力を $x$ と $y$ の成分に分けて、1つずつ考える}
\noindent\begin{minipage}[c]{0.46\linewidth}
\centering
\begin{tikzpicture}[x=1mm, y=1mm]
  \draw[bkGray, line width=0.6pt, -{Stealth[length=4pt]}] (-4,0) -- (52,0) node[anchor=west, font=\sffamily\small, text=bkInk] {$x$};
  \draw[bkGray, line width=0.6pt, -{Stealth[length=4pt]}] (0,-4) -- (0,44) node[anchor=south, font=\sffamily\small, text=bkInk] {$y$};
  \draw[bkOrange!70!black, line width=2pt] (0,0) -- (30,20);
  \draw[bkGray, line width=0.4pt, dashed] (30,0) -- (30,20) -- (0,20);
  \fill[bkInk] (30,20) circle (1);
  \draw[bk arrow, bkGreen!70!black, line width=1.6pt] (30,20) -- (30,38);
  \draw[bk arrow, bkBlue, line width=1.6pt] (30,20) -- (44,20);
  \node[font=\sffamily\footnotesize, anchor=north west, inner sep=1pt] at (30.8,19.4) {P$(x,\ y)$};
  \node[font=\sffamily\footnotesize, text=bkGreen!50!black, anchor=west, inner sep=1.5pt] at (30.6,32) {$F_y$};
  \node[font=\sffamily\footnotesize, text=bkBlue!80!black, anchor=south, inner sep=1.5pt] at (40,20.4) {$F_x$};
  \node[font=\sffamily\footnotesize, anchor=north, inner sep=1.5pt] at (30,-0.6) {$x$};
  \node[font=\sffamily\footnotesize, anchor=east, inner sep=1.5pt] at (-0.6,20) {$y$};
  \node[font=\sffamily\footnotesize, anchor=north east, inner sep=1pt] at (-0.5,-0.5) {O};
\end{tikzpicture}
\end{minipage}\hfill
\begin{minipage}[c]{0.5\linewidth}
\textbf{上向きの $F_y$（緑）}\par
上向きの力のうでの長さは、Oから横に $x$。Pを上に押すと、Oのまわりを左回りに回す。だから $+xF_y$。\par\vspace{4pt}
\textbf{右向きの $F_x$（青）}\par
右向きの力のうでの長さは、Oからの高さ $y$。Pを右に押すと、Oのまわりを右回り（時計回り）に回す。だから $-yF_x$。
\end{minipage}

\par\vspace{4pt}
2つを合わせて、
\bkivoline{\fbox{$M=xF_y-yF_x$}\qquad（＋なら左回り、－なら右回り）}
$x$ や $y$ がマイナスの点、$F_x$ や $F_y$ がマイナスの力でも、この式のままで、符号まで正しく出てきます。

\bkivostep{3}{数で、2通りにたしかめる}
P$(0.3,\ 0.4)$（m）を、真上向きに\qty{10}{N}で押します。$\vec F=(0,\ 10)$ です。
\begin{itemize}[leftmargin=1.2\zw, itemsep=2pt, topsep=2pt, label={\color{\bkTT}●}]
\item \textbf{成分の式で}\quad $M=0.3\times10-0.4\times0=3$（\unit[inter-unit-product=\ensuremath{\cdot}]{N.m}、左回り）
\item \textbf{②の式で}\quad うでの長さは $r=\sqrt{0.3^2+0.4^2}=0.5$（4-4）。うでの向きは $x$ 軸から約53.1°（$\cos$ が $\dfrac{0.3}{0.5}=0.6$ の角）、力は90°の向きなので、なす角は $\theta=90°-53.1°=36.9°$。$\sin36.9°\approx0.6$ なので、$M=0.5\times10\times0.6=3$
\end{itemize}
同じ3になりました。成分の式なら、角度を求めずにすみます。

\bksubsection{4}{回す向きを、矢印で表す：右ねじ}

\leavevmode\label{bk:4-14:migineji}平面の上なら、回る向きは左回りか右回りの2つだけで、符号で区別できました。でも空間では、回る「軸」がいろいろな向きを向きます。ドアは上下の軸のまわりに、自転車の車輪は横の軸のまわりに回ります。回る向きを、どう表せばよいでしょう。

\noindent\begin{minipage}[c]{0.56\linewidth}
答えは、\textbf{回る軸の向きの矢印}で表すことです。軸には2つの向きがあるので、どちらを向けるかは、右ねじで決めます（4-13）。
\begin{itemize}[leftmargin=1.2\zw, itemsep=2pt, topsep=2pt, label={\color{\bkTT}●}]
\item 右手の4本の指を、回る向きにそって曲げる。
\item そのとき親指が向く向きが、矢印の向き。
\end{itemize}
右ねじを、回る向きに回したとき、ねじが進む向きと同じです。ペットボトルのふたを上から見て右回り（時計回り）に回すと、ふたは下へ進み、しまります。この回り方は、下向きの矢印で表します。
\end{minipage}\hfill
\begin{minipage}[c]{0.4\linewidth}
\centering
\begin{tikzpicture}[x=1mm, y=1mm]
  % ペットボトルのふた（上から見ると右回り）
  \fill[bkBlue!15] (-9,-14) rectangle (9,4);
  \draw[bkBlue!60!black, line width=0.5pt] (-9,-14) -- (-9,4) (9,-14) -- (9,4);
  \fill[bkOrange!45] (-7,4) rectangle (7,11);
  \draw[bkOrange!70!black, line width=0.6pt] (-7,4) rectangle (7,11);
  \foreach \k in {-5,-3,...,5} {\draw[bkOrange!70!black, line width=0.4pt] (\k,4.5) -- (\k,10.5);}
  \fill[bkOrange!25] (0,11) ellipse (7 and 2);
  \draw[bkOrange!70!black, line width=0.6pt] (0,11) ellipse (7 and 2);
  \draw[bkRed, line width=1.3pt, -{Stealth[length=5pt]}] (-11,15.5) arc (160:20:11.7 and 3.4);
  \node[font=\sffamily\footnotesize, text=bkRed!80!black, anchor=south] at (0,19.5) {上から見て右回り};
  \draw[bk arrow, bkGreen!70!black, line width=2pt] (17,8) -- (17,-12);
  \node[font=\sffamily\footnotesize, text=bkGreen!50!black, anchor=west, align=left] at (18.5,-2) {ふたの\\進む向き};
\end{tikzpicture}
\end{minipage}

\par\vspace{6pt}
\bkivohead{紙の上での書き方：$\odot$ と $\otimes$}
平面の図で、左回り・右回りを矢印で表すと、矢印は紙に垂直になります。そこで、紙から手前へ出てくる向きを $\odot$、紙のおくへ入る向きを $\otimes$ と書きます。矢の先を前から見た形が $\odot$、矢の羽根を後ろから見た形が $\otimes$ です。

\begin{bkfigure}
\begin{tikzpicture}[x=1mm, y=1mm]
  \draw[bkGray, line width=0.6pt] (0,0) circle (12);
  \draw[bkRed, line width=1.3pt, -{Stealth[length=5pt]}] (12,0) arc (0:300:12);
  \bkivoout{(0,0)}{3.2mm}
  \node[font=\sffamily\small, anchor=north] at (0,-15) {左回り → $\odot$（手前へ）};
  \begin{scope}[shift={(54,0)}]
  \draw[bkGray, line width=0.6pt] (0,0) circle (12);
  \draw[bkRed, line width=1.3pt, -{Stealth[length=5pt]}] (12,0) arc (0:-300:12);
  \bkivoin{(0,0)}{3.2mm}
  \node[font=\sffamily\small, anchor=north] at (0,-15) {右回り → $\otimes$（おくへ）};
  \end{scope}
  % 矢の先・羽根
  \begin{scope}[shift={(90,2)}]
  \draw[bkInk, line width=1pt] (0,0) -- (16,0);
  \fill[bkInk] (16,-2) -- (20,0) -- (16,2) -- cycle;
  \draw[bkInk, line width=0.8pt] (0,0) -- (-3,3) (0,0) -- (-3,-3) (2,0) -- (-1,3) (2,0) -- (-1,-3);
  \node[font=\sffamily\footnotesize, anchor=north, align=center] at (10,-5) {先を見ると $\odot$\\羽根を見ると $\otimes$};
  \end{scope}
\end{tikzpicture}
\bkfigcaption{$xy$ 平面（紙）の上で左回りに回すと、右ねじは紙から手前（$z$ の＋の向き）へ進む。4-13の右手系の約束どおり。}
\end{bkfigure}

$xy$ 平面で左回り（$x$ から $y$ へ回る向き）は、4-13の右手系の約束で、$z$ 軸の＋の向き（$\odot$）です。だから、③の $M=xF_y-yF_x$ は、「回す向きの矢印の $z$ 成分」と読めます。$M$ が＋なら矢印は $\odot$、－なら $\otimes$ です。

\bksubsection{5}{2つのベクトルから、垂直なベクトルを作る：外積}

\leavevmode\label{bk:4-14:gaiseki}ここまでで、回す働きは「うでのベクトル」と「力のベクトル」から決まり、大きさは $rF\sin\theta$、向きは右ねじの軸の向き、とわかりました。これを、どんな2つのベクトルにも使えるように、計算のきまりにします。

\bkivostep{1}{外積の決め方}
2つのベクトル $\vec a$、$\vec b$ のなす角を $\theta$ とします。次の大きさと向きをもつベクトルを、$\vec a$ と $\vec b$ の\textbf{外積}（がいせき）といい、$\vec a\times\vec b$ と書きます。\bkidx{がいせき}{外積}
\begin{bkcenter}
\begin{bktable}{colspec={Q[l,m,wd=22mm] Q[l,m,wd=94mm]}, rowsep=2pt}
大きさ & $|\vec a\times\vec b|=|\vec a||\vec b|\sin\theta$ \\
向き & $\vec a$ にも $\vec b$ にも垂直。$\vec a$ を $\vec b$ のほうへ（小さいほうの角 $\theta$ だけ）回すとき、右ねじが進む向き。 \\
\end{bktable}
\end{bkcenter}
「×」の記号は、ふつうの数のかけ算と同じ形ですが、ベクトルの外積ではベクトルが答えです。読み方は「エー・クロス・ビー」。内積の「・」と取りちがえないようにします。

\noindent\begin{minipage}[c]{0.5\linewidth}
\centering
\begin{tikzpicture}[ivo3d=9.5mm]
  \filldraw[fill=bkTheme!25, draw=bkDeep, line width=0.6pt] (0,0,0) -- (2.4,0,0) -- (2.4,2.6,0) -- (0,2.6,0) -- cycle;
  \fill[bkTheme!45] (0,0,0) -- (2.4,0,0) -- (0,2.6,0) -- cycle;
  \draw[bk arrow, bkBlue, line width=1.6pt] (0,0,0) -- (2.4,0,0);
  \draw[bk arrow, bkOrange, line width=1.6pt] (0,0,0) -- (0,2.6,0);
  \draw[bk arrow, bkRed, line width=2pt] (0,0,0) -- (0,0,2.6);
  \draw[bkInk, line width=1pt, -{Stealth[length=4pt]}] (0.9,0,0) .. controls (0.9,0.6,0) and (0.6,0.9,0) .. (0,0.9,0);
  \node[font=\sffamily\small, text=bkBlue!80!black, anchor=east, inner sep=2pt] at (2.4,0,0) {$\vec a$};
  \node[font=\sffamily\small, text=bkOrange!80!black, anchor=west, inner sep=2pt] at (0,2.6,0) {$\vec b$};
  \node[font=\sffamily\small, text=bkRed!80!black, anchor=south, inner sep=2pt] at (0,0,2.6) {$\vec a\times\vec b$};
  \node[font=\sffamily\footnotesize, anchor=north west, inner sep=1pt] at ($(1.2,1.3,0)+(1mm,0)$) {面積};
\end{tikzpicture}
\end{minipage}\hfill
\begin{minipage}[c]{0.46\linewidth}
$\vec a$ と $\vec b$ がつくる平行四辺形（黄緑）を考えます。\par\vspace{3pt}
\textbf{大きさ}\quad 平行四辺形の面積は、三角形（こい黄緑）の2倍です。三角形の面積は $\dfrac{1}{2}|\vec a||\vec b|\sin\theta$（4-7）なので、平行四辺形は $|\vec a||\vec b|\sin\theta$。外積の大きさは、\textbf{平行四辺形の面積}と同じです。\par\vspace{3pt}
\textbf{向き}\quad 平行四辺形の面に垂直（面の法線ベクトル、4-13）。
\end{minipage}

\par\vspace{4pt}
力のモーメントは、うで $\vec r$ と力 $\vec F$ の外積 $\vec r\times\vec F$ です。大きさ $rF\sin\theta$ は②の式、向きは④の右ねじの向きで、ぴったり合っています。

\bkivostep{2}{基本ベクトルどうしの外積}
4-13の基本ベクトル $\vec i$、$\vec j$、$\vec k$ で外積を計算します。どれも長さ1で、たがいに直角（$\sin90°=1$）なので、外積の大きさは1です。向きは右ねじで決まります。
\begin{itemize}[leftmargin=1.2\zw, itemsep=2pt, topsep=2pt, label={\color{\bkTT}●}]
\item $\vec i\times\vec j$：$x$ から $y$ へ回すと、右ねじは $z$ へ進む（右手系の約束）。だから $\vec i\times\vec j=\vec k$。
\item 同じように、$y$ から $z$ へ回すと $x$ へ、$z$ から $x$ へ回すと $y$ へ進む。だから $\vec j\times\vec k=\vec i$、$\vec k\times\vec i=\vec j$。
\item 順番を逆にすると、回す向きが逆なので、ねじは反対へ進む。$\vec j\times\vec i=-\vec k$ など。
\item 同じものどうしは、なす角が0°で $\sin0°=0$。$\vec i\times\vec i=\vec 0$（ゼロベクトル。4-10）。
\end{itemize}

\begin{bkfigure}
\begin{tikzpicture}[x=1mm, y=1mm]
  \node[circle, draw=bkDeep, line width=0.9pt, fill=bkTheme!12, minimum size=10mm, font=\large] (i) at (90:13) {$\vec i$};
  \node[circle, draw=bkDeep, line width=0.9pt, fill=bkTheme!12, minimum size=10mm, font=\large] (j) at (-30:13) {$\vec j$};
  \node[circle, draw=bkDeep, line width=0.9pt, fill=bkTheme!12, minimum size=10mm, font=\large] (k) at (210:13) {$\vec k$};
  \draw[bk link, bkRed, line width=1.1pt] (i) to[bend left=25] (j);
  \draw[bk link, bkRed, line width=1.1pt] (j) to[bend left=25] (k);
  \draw[bk link, bkRed, line width=1.1pt] (k) to[bend left=25] (i);
  \node[font=\sffamily\small, align=left, anchor=west] at (30,0) {矢印の向きにたどると＋：\\[2pt]$\vec i\times\vec j=\vec k$、\enspace$\vec j\times\vec k=\vec i$、\enspace$\vec k\times\vec i=\vec j$\\[4pt]逆向きにたどると－：\\[2pt]$\vec j\times\vec i=-\vec k$、\enspace$\vec k\times\vec j=-\vec i$、\enspace$\vec i\times\vec k=-\vec j$};
\end{tikzpicture}
\bkfigcaption{$x\to y\to z\to x$ の順にまわる。この順の2つをかけると、残りの1つ（＋）。}
\end{bkfigure}

\bkivostep{3}{成分の式を作る}
$\vec a=a_x\vec i+a_y\vec j+a_z\vec k$、$\vec b=b_x\vec i+b_y\vec j+b_z\vec k$ と書きます（4-13）。外積でも、かっこをはずすきまり（分配法則。$(\vec a+\vec b)\times\vec c=\vec a\times\vec c+\vec b\times\vec c$ など）と、数を前に出すきまりが成り立つことが知られています。それを使って、全部で9組のかけ算をします。同じものどうし（$\vec i\times\vec i$ など）の3組は0なので、残りは6組です。
\begin{align*}
\vec a\times\vec b&=a_xb_y\,(\vec i\times\vec j)+a_xb_z\,(\vec i\times\vec k)+a_yb_x\,(\vec j\times\vec i)\\
&\quad+a_yb_z\,(\vec j\times\vec k)+a_zb_x\,(\vec k\times\vec i)+a_zb_y\,(\vec k\times\vec j)\\
&=a_xb_y\,\vec k-a_xb_z\,\vec j-a_yb_x\,\vec k+a_yb_z\,\vec i+a_zb_x\,\vec j-a_zb_y\,\vec i
\end{align*}
$\vec i$、$\vec j$、$\vec k$ ごとにまとめると、
\bkivoline{\fbox{$\vec a\times\vec b=(a_yb_z-a_zb_y,\ \ a_zb_x-a_xb_z,\ \ a_xb_y-a_yb_x)$}}
$z$ 成分の $a_xb_y-a_yb_x$ は、③の平面の式 $xF_y-yF_x$ と同じ形です。$x$ 成分と $y$ 成分は、文字を $x\to y\to z\to x$ と1つずつずらすと作れます。

\bkivohead{たすきがけで、まちがえずに計算する}
\begin{bkfigure}
\begin{tikzpicture}[x=1mm, y=1mm]
  \foreach \c/\t in {0/x,1/y,2/z,3/x,4/y} {
    \node[font=\normalsize] (a\c) at (\c*14,8) {$a_\t$};
    \node[font=\normalsize] (b\c) at (\c*14,-4) {$b_\t$};}
  \draw[bkGray, line width=0.4pt, dashed] (35,13) -- (35,-9);
  \draw[bkBlue, line width=1.1pt] (a1.south east) -- (b2.north west);
  \draw[bkRed, line width=1.1pt] (a2.south west) -- (b1.north east);
  \node[font=\sffamily\footnotesize, text=bkGray, anchor=south] at (49,13) {くり返し};
  \node[font=\sffamily\small, align=left, anchor=west] at (68,2) {$x$ 成分：\textcolor{bkBlue}{$a_yb_z$}$-$\textcolor{bkRed!80!black}{$a_zb_y$}\\[3pt]$y$ 成分：右へ1つずらして $a_zb_x-a_xb_z$\\[3pt]$z$ 成分：もう1つずらして $a_xb_y-a_yb_x$};
\end{tikzpicture}
\bkfigcaption{成分を2回くり返して並べる。$x$ 成分は、$y$ と $z$ の列でたすきにかけ、右下がり（青）－左下がり（赤）。}
\end{bkfigure}

\bkivostep{4}{数で、たしかめる}
$\vec a=(3,\ 0,\ 0)$、$\vec b=(1,\ 2,\ 0)$ では、
\[ \vec a\times\vec b=(0\times0-0\times2,\ \ 0\times1-3\times0,\ \ 3\times2-0\times1)=(0,\ 0,\ 6) \]
\begin{itemize}[leftmargin=1.2\zw, itemsep=2pt, topsep=2pt, label={\color{\bkTT}●}]
\item \textbf{向き}\quad $\vec a$、$\vec b$ は床（$xy$ 平面）の上の矢印で、答えは真上（$z$ の＋）の向き。$\vec a$ から $\vec b$ へは左回りなので、右ねじは上へ進む。合っています。
\item \textbf{大きさ}\quad 6。平行四辺形の底辺は $|\vec a|=3$、高さは $\vec b$ の $y$ 成分の2なので、面積は $3\times2=6$。合っています。
\end{itemize}
もう1つ、$\vec a=(1,\ 2,\ 0)$、$\vec b=(0,\ 1,\ 3)$ では、$\vec a\times\vec b=(6,\ -3,\ 1)$ です（例題4）。$\vec a$ との内積は $6-6+0=0$、$\vec b$ との内積は $0-3+3=0$。どちらとも内積が0なので、\textbf{答えは $\vec a$ にも $\vec b$ にも垂直}です（4-12）。

\bkivostep{5}{表計算で、外積を計算する}
成分の式は、表計算にそのまま入れられます。2行めに $\vec a$、3行めに $\vec b$ の成分を入れ、4行めに外積の式を入れます。
\begin{bkcenter}
\bksheet{4}{26}{%
1 & & $x$ 成分 & $y$ 成分 & $z$ 成分 \\
2 & $\vec a$ & 1 & 2 & 0 \\
3 & $\vec b$ & 0 & 1 & 3 \\
4 & $\vec a\times\vec b$ & \texttt{=C2*D3-D2*C3} & \texttt{=D2*B3-B2*D3} & \texttt{=B2*C3-C2*B3} \\
5 & （結果） & 6 & $-3$ & 1 \\}
\bktabcaption{4行めの式は、たすきがけ（③）をセルの名前で書いたもの。5行めは、4行めに出る答え。}
\end{bkcenter}
2行め・3行めの数を変えるだけで、どんな外積でもすぐに計算できます。例題や問題の答えのたしかめに使いましょう。

\bksubsection{6}{外積の性質：内積とくらべる}

\bkivostep{1}{順番を入れかえると、向きが反対}\label{bk:4-14:seishitsu}
$\vec b\times\vec a$ は、$\vec b$ から $\vec a$ へ回すので、右ねじは $\vec a\times\vec b$ と反対へ進みます。大きさは同じです。
\bkivoline{\fbox{$\vec b\times\vec a=-\,\vec a\times\vec b$}}
ふつうの数のかけ算は $2\times3=3\times2$ ですが、外積では順番が大切です。成分の式でも、$a$ と $b$ を入れかえると、どの成分も符号が反対になります。

\bkivostep{2}{平行なら0}
$\vec a$ と $\vec b$ が同じ向き（$\theta=0°$）か反対向き（$\theta=180°$）なら、$\sin\theta=0$ なので、外積は $\vec 0$ です。平行四辺形がつぶれて、面積が0になるからです。とくに $\vec a\times\vec a=\vec 0$、$\vec a\times(2\vec a)=\vec 0$。ドアのノブを、ちょうつがいに向かって押しても回らないのと同じです。

\begin{tsumazuki}{外積の順番を入れかえる}
\bkng{$\vec j\times\vec i=\vec k$}
\bkok{$\vec j\times\vec i=-\vec k$}
外積では順番が大切です。$x\to y\to z\to x$ の順なら＋、逆なら－です。$\vec b\times\vec a=-\vec a\times\vec b$。
\end{tsumazuki}

\begin{bkcenter}
\begin{bktable}{colspec={Q[l,m,wd=30mm] Q[c,m,wd=42mm] Q[c,m,wd=44mm]}, rowsep=2pt, row{1}={font=\sffamily\small}}
 & 内積 $\vec a\cdot\vec b$（4-12） & 外積 $\vec a\times\vec b$ \\
答え & ふつうの数（スカラー） & ベクトル \\
大きさ & $|\vec a||\vec b|\cos\theta$ & $|\vec a||\vec b|\sin\theta$ \\
いちばん大きいとき & 平行（$\theta=0°$） & 直角（$\theta=90°$） \\
0になるとき & 直角 & 平行 \\
順番を入れかえると & 同じ & 符号が反対 \\
物理では & 仕事 $W=\vec F\cdot\vec s$ & 力のモーメント $\vec r\times\vec F$ \\
\end{bktable}
\bktabcaption{内積は「どれだけ同じ向きか」、外積は「どれだけ直角に開いているか」を表す。}
\end{bkcenter}

\begin{motto}[4]
\textbf{内積と外積の、うれしい関係}\par
$\cos^2\theta+\sin^2\theta=1$（4-8）の両辺に $|\vec a|^2|\vec b|^2$ をかけると、
\[ (|\vec a||\vec b|\cos\theta)^2+(|\vec a||\vec b|\sin\theta)^2=|\vec a|^2|\vec b|^2 \]
1つめは内積の2乗、2つめは外積の大きさの2乗です。つまり
\[ (\vec a\cdot\vec b)^2+|\vec a\times\vec b|^2=|\vec a|^2|\vec b|^2 \]
⑤の $\vec a=(1,\ 2,\ 0)$、$\vec b=(0,\ 1,\ 3)$ でたしかめます。$\vec a\cdot\vec b=2$、$\vec a\times\vec b=(6,\ -3,\ 1)$ なので、左は $2^2+(36+9+1)=4+46=50$。右は $|\vec a|^2|\vec b|^2=5\times10=50$。合っています。「同じ向きの部分」と「直角に開いた部分」を合わせると、全体になるのです。
\end{motto}

\bkivostep{3}{三角形の面積を、座標から求める}
3点A、B、Cがつくる三角形の面積は、$\overrightarrow{\mathrm{AB}}$ と $\overrightarrow{\mathrm{AC}}$ がつくる平行四辺形の半分です。だから
\bkivoline{\fbox{三角形ABCの面積 $=\dfrac{1}{2}\,\bigl|\overrightarrow{\mathrm{AB}}\times\overrightarrow{\mathrm{AC}}\bigr|$}}
4-7の $\dfrac{1}{2}ab\sin C$ と同じことを、角度を求めずに、座標だけで計算できます。空間の中の、ななめにかたむいた三角形でも使えます。

\bkivostep{4}{3点を通る平面の式}
4-13では、法線ベクトル $(a,\ b,\ c)$ がわかれば、平面の式 $ax+by+cz=d$ が作れました。外積を使うと、3点を通る平面の法線ベクトルが作れます。$\overrightarrow{\mathrm{AB}}\times\overrightarrow{\mathrm{AC}}$ は、$\overrightarrow{\mathrm{AB}}$ にも $\overrightarrow{\mathrm{AC}}$ にも垂直、つまり平面ABCに垂直だからです。

A$(1,\ 0,\ 0)$、B$(0,\ 2,\ 0)$、C$(0,\ 0,\ 2)$ では、$\overrightarrow{\mathrm{AB}}=(-1,\ 2,\ 0)$、$\overrightarrow{\mathrm{AC}}=(-1,\ 0,\ 2)$ なので、成分の式で $\overrightarrow{\mathrm{AB}}\times\overrightarrow{\mathrm{AC}}=(4,\ 2,\ 2)$ です（くわしい計算は例題5）。だから平面の式は $4x+2y+2z=d$。Aを入れると $4\times1+0+0=4$ なので $d=4$。両辺を2でわって
\[ 2x+y+z=2 \]
Bを入れると $0+2+0=2$、Cを入れると $0+0+2=2$ で、3点とも式をみたしています。

\bksubsection{7}{トルクと角運動量}

\bkivostep{1}{空間のトルク $\vec r\times\vec F$}\label{bk:4-14:torque}
回転の中心から力をかける点までのベクトルを $\vec r$、力を $\vec F$ とすると、力のモーメント（トルク）は
\bkivoline{\fbox{$\vec N=\vec r\times\vec F$}\qquad（大きさ $rF\sin\theta$、向きは回す軸の向き）}
です。トルクは $\vec N$ や、ギリシャ文字の $\tau$（タウ）で書くことが多いです。

\noindent\begin{minipage}[c]{0.5\linewidth}
\centering
\begin{tikzpicture}[ivo3d=11mm]
  \bkivoaxes{2.4}{3.6}{2.2}
  \draw[bk arrow, bkGreen!70!black, line width=2pt] (0,0,0) -- (-1.6,0,0);
  \node[font=\sffamily\footnotesize, text=bkGreen!50!black, anchor=south west, inner sep=1.5pt] at (-1.6,0,0) {$\vec N$（おくへ）};
  \fill[bkGray!55] (0,0) circle (3.2mm);
  \draw[bkGray!80!black, line width=0.6pt] (0,0) circle (3.2mm);
  \fill[bkGray!80!black] (0,0) circle (1mm);
  \draw[bkOrange!70!black, line width=4pt, line cap=round] (0,0.35,0) -- (0,2.5,0);
  \draw[bk arrow, bkRed, line width=1.6pt] (0,2.5,1.5) -- (0,2.5,0.18);
  \node[font=\sffamily\footnotesize, text=bkRed!80!black, anchor=west, inner sep=2pt] at (0,2.5,1.1) {$\vec F$（下へ\qty{40}{N}）};
  \node[font=\sffamily\footnotesize, anchor=north, inner sep=2pt] at (0,1.5,-0.1) {スパナ \qty{0.25}{m}};
  \node[font=\sffamily\footnotesize, anchor=north east, inner sep=1pt] at (0,-0.3,-0.25) {ボルト};
\end{tikzpicture}
\end{minipage}\hfill
\begin{minipage}[c]{0.46\linewidth}
かべにさしたボルト（軸は $x$ の向き。$x$ の＋がこちら向き）に、$y$ の向きにスパナ（長さ\qty{0.25}{m}）をかけ、先を真下に\qty{40}{N}で押します。$\vec r=(0,\ 0.25,\ 0)$、$\vec F=(0,\ 0,\ -40)$ です。
\end{minipage}

\par\vspace{4pt}
外積の成分の式（⑤）で計算します。
\begin{itemize}[leftmargin=1.2\zw, itemsep=1pt, topsep=2pt, label={\color{\bkTT}●}]
\item $x$ 成分\enspace $0.25\times(-40)-0\times0=-10$
\item $y$ 成分\enspace $0\times0-0\times(-40)=0$
\item $z$ 成分\enspace $0\times0-0.25\times0=0$
\end{itemize}
トルクは $\vec N=(-10,\ 0,\ 0)$ で、$x$ の－の向き（かべのおくへ）に\bkivoNm{10}。手前から見るとスパナは右回り（時計回り）に回り、右ねじはおくへ進みます。ボルトがかべにしめこまれる向きです。
\par\vspace{4pt}
トルクの矢印の向きが、ボルトの軸の向きと同じになりました。トルクの矢印は、「どの軸のまわりに、どちら向きに回すか」を1本の矢印で表しているのです。

\bkivostep{2}{角運動量：回る勢い}
質量 $m$ の物が、速度 $\vec v$ で動いています。$m\vec v$ を運動量といいます（物理シリーズで学びます）。中心から物までのベクトルを $\vec r$ として、
\bkivoline{\fbox{$\vec L=\vec r\times m\vec v$}}
を、中心のまわりの\textbf{角運動量}といいます。\bkidx{かくうんどうりょう}{角運動量}「どの軸のまわりに、どれだけの勢いで回っているか」を表すベクトルです。

半径 $r$ の円の上を速さ $v$ で回る物では、速度は半径に直角（4-2の接線）なので、$\sin90°=1$ で、大きさは $L=rmv$ です。$v=r\omega$（4-8）を入れると $L=mr^2\omega$。

フィギュアスケートの選手は、広げたうでをたたむと、速く回ります。外から回す力（トルク）がはたらかなければ、角運動量は変わらない、という自然の法則があるからです（角運動量の保存。物理シリーズ 力-6）。うでの物の回る半径 $r$ が半分になると、$r^2$ は $\dfrac{1}{4}$ になるので、$mr^2\omega$ が変わらないためには、$\omega$ が4倍になります（体のほかの部分は考えない、簡単にしたモデル）。
\begin{bkfigure}
\begin{tikzpicture}[x=1mm, y=1mm]
  \fill[bkBlue!20] (0,0) circle (5);
  \draw[bkBlue!70!black, line width=0.6pt] (0,0) circle (5);
  \draw[bkInk!70, line width=1.4pt] (-24,0) -- (-5,0) (5,0) -- (24,0);
  \fill[bkOrange] (-24,0) circle (1.6); \fill[bkOrange] (24,0) circle (1.6);
  \draw[bkRed, line width=1.2pt, -{Stealth[length=4pt]}] (24,2) -- (24,9.5);
  \draw[bkRed, line width=1.2pt, -{Stealth[length=4pt]}] (-24,-2) -- (-24,-9.5);
  \node[font=\sffamily\small, anchor=north] at (0,-19) {うでを広げる：$r$、$\omega$};
  \begin{scope}[shift={(72,0)}]
  \fill[bkBlue!20] (0,0) circle (5);
  \draw[bkBlue!70!black, line width=0.6pt] (0,0) circle (5);
  \draw[bkInk!70, line width=1.4pt] (-12,0) -- (-5,0) (5,0) -- (12,0);
  \fill[bkOrange] (-12,0) circle (1.6); \fill[bkOrange] (12,0) circle (1.6);
  \draw[bkRed, line width=1.2pt, -{Stealth[length=4pt]}] (12,2) -- (12,17);
  \draw[bkRed, line width=1.2pt, -{Stealth[length=4pt]}] (-12,-2) -- (-12,-17);
  \node[font=\sffamily\small, anchor=north] at (0,-19) {うでをたたむ：$\dfrac{r}{2}$、$4\omega$};
  \end{scope}
  \draw[bk arrow, bkGray, line width=1pt] (34,0) -- (52,0);
\end{tikzpicture}
\bkfigcaption{上から見たスケートの選手（うでの先のおもりだけを考えた、簡単なモデル）。回る半径が半分になると、角速度は4倍。赤い矢印はおもりの速さ $v=r\omega$ で、$\dfrac{r}{2}\times4\omega=2r\omega$ なので2倍になる。}
\end{bkfigure}

\bksubsection{8}{電流が磁場から受ける力}

\leavevmode\label{bk:4-14:jiba}モーターが回るのは、磁石の近くの導線に電流を流すと、導線が力を受けるからです。この力の向きも、外積で表せます。

\bkivostep{1}{力の式}
磁石の力が届いている場所を磁場といい、その向きと強さをベクトル $\vec B$（単位 T、テスラ）で表します。電荷 $q$（単位 C、クーロン）の粒子が速度 $\vec v$ で磁場の中を動くと、
\bkivoline{\fbox{$\vec F=q\,\vec v\times\vec B$}}
の力を受けます。この力を\textbf{ローレンツ力}といいます。\bkidx{ろーれんつりょく}{ローレンツ力}導線の中の電流は、たくさんの電荷の流れです。長さ $L$ の導線に電流 $I$ が流れるとき、電流の向きで長さ $L$ のベクトルを $\vec L$ とすると、導線が受ける力は
\bkivoline{\fbox{$\vec F=I\,\vec L\times\vec B$}\qquad（大きさ $ILB\sin\theta$）}

\noindent\begin{minipage}[c]{0.48\linewidth}
\centering
\begin{tikzpicture}[ivo3d=10mm]
  \bkivoaxes{3.4}{3.6}{2.4}
  % 磁石（N極とS極）
  \fill[bkRed!55] (1.5,-1.2,0.2) ++(-3mm,-4mm) rectangle ++(6mm,8mm);
  \node[font=\sffamily\small\bfseries, text=bkInk] at (1.5,-1.2,0.2) {N};
  \fill[bkBlue!45] (1.5,3.2,0.2) ++(-3mm,-4mm) rectangle ++(6mm,8mm);
  \node[font=\sffamily\small\bfseries, text=bkInk] at (1.5,3.2,0.2) {S};
  \draw[bkGray, line width=0.5pt, dashed, -{Stealth[length=3pt]}] (1.5,-0.7,0.9) -- (1.5,2.7,0.9);
  \node[font=\sffamily\footnotesize, text=bkInk, anchor=south, inner sep=1.5pt] at (1.5,-0.2,0.9) {$\vec B$};
  % 導線（x の向き）
  \draw[bkOrange!70!black, line width=2.4pt] (3,1,0.2) -- (0,1,0.2);
  \draw[bk arrow, bkOrange!80!black, line width=1.2pt] (0.4,1,0.2) -- (2.6,1,0.2);
  \node[font=\sffamily\footnotesize, text=bkOrange!80!black, anchor=north, inner sep=2pt] at (3,1,0.2) {電流 $I$};
  \draw[bk arrow, bkRed, line width=2pt] (1.5,1,0.35) -- (1.5,1,1.8);
  \node[font=\sffamily\footnotesize, text=bkRed!80!black, anchor=west, inner sep=2pt] at (1.5,1,1.7) {力 $\vec F$};
\end{tikzpicture}
\end{minipage}\hfill
\begin{minipage}[c]{0.48\linewidth}
電流が $x$ の＋の向き、磁場が $y$ の＋の向きなら、$\vec L$ は $\vec i$ の向き、$\vec B$ は $\vec j$ の向きです。$\vec i\times\vec j=\vec k$ なので、導線は $z$ の＋の向き（上）に押されます。

$I=\qty{3}{A}$、$L=\qty{0.2}{m}$、$B=\qty{0.5}{T}$ なら、直角なので $\sin90°=1$ で、
\[ F=3\times0.2\times0.5=\qty{0.3}{N} \]
\end{minipage}

\bkivostep{2}{フレミングの左手の法則}
中学校の理科では、この向きを左手で覚えました。左手の親指・人さし指・中指を直角に開き、
\par\vspace{2pt}\noindent\begin{minipage}[c]{0.56\linewidth}
\small
\begin{bktable}{colspec={Q[l,m,wd=17mm] Q[l,m,wd=23mm] Q[l,m,wd=25mm]}, rowsep=1.5pt}
左手の指 & 表すもの & 外積では \\
中指 & 電流 $I$（$\vec L$） & 外積の1つめ \\
人さし指 & 磁場 $\vec B$ & 外積の2つめ \\
親指 & 力 $\vec F$ & 答えの向き \\
\end{bktable}
\end{minipage}\hfill
\begin{minipage}[c]{0.4\linewidth}
\centering
\begin{tikzpicture}[ivo3d=10mm]
  \fill[bkOrange!22] (0,0,0) ++(-4mm,-4.5mm) rectangle ++(8mm,7mm);
  \draw[bkOrange!70!black, line width=0.6pt, rounded corners=2mm] (0,0,0) ++(-4mm,-4.5mm) rectangle ++(8mm,7mm);
  \bkivofinger{(0,0,0)}{(2.0,0,0)}{3.2mm}
  \bkivofinger{(0,0,0)}{(0,2.2,0)}{3.2mm}
  \bkivofinger{(0,0,0)}{(0,0,2.1)}{3.4mm}
  \node[font=\sffamily\footnotesize, anchor=east, inner sep=2pt, align=right] at ($(1.9,0,0)+(-2mm,0)$) {中指\\電流};
  \node[font=\sffamily\footnotesize, anchor=north, inner sep=2pt] at ($(0,1.4,0)+(0,-2.4mm)$) {人さし指：磁場};
  \node[font=\sffamily\footnotesize, anchor=west, inner sep=2pt] at ($(0,0,1.6)+(2.4mm,0)$) {親指：力};
\end{tikzpicture}
\end{minipage}\par\vspace{4pt}
「電・磁・力」と、中指から順に覚えます。外積 $\vec L\times\vec B$ の右ねじの向きと、左手の親指の向きは、いつも同じになります。左手の法則は、外積の向きを、右ねじを使わずに指で調べる方法なのです。外積のきまりを知っていれば、左手の法則を覚えていなくても、向きを計算で出せます。

\bkivostep{3}{マイナスの電荷のとき}
電子の電荷は負（$q=-1.6\times10^{-19}$\,C）です。$q$ がマイナスなので、$\vec F=q\,\vec v\times\vec B$ は、$\vec v\times\vec B$ と反対向きになります。電子が $x$ の＋の向きに動いても、電流は電子と反対向き（$x$ の－の向き）と決められているので、電流で考えても同じ答えになります。

\begin{motto}[5]
\textbf{3つのベクトルでつくる箱の体積}\par
$\vec a$、$\vec b$、$\vec c$ の3本がつくる、ななめにかたむいた箱（平行六面体）の体積を考えます。底の平行四辺形の面積は $|\vec a\times\vec b|$ です。高さは、$\vec c$ の、底に垂直な向き（$\vec a\times\vec b$ の向き）の成分です（4-12の正射影）。だから
\[ \text{体積}=(\vec a\times\vec b)\cdot\vec c \]
たとえば $\vec a=(2,\ 0,\ 0)$、$\vec b=(0,\ 3,\ 0)$、$\vec c=(1,\ 1,\ 4)$ なら、$\vec a\times\vec b=(0,\ 0,\ 6)$ で、$(0,\ 0,\ 6)\cdot(1,\ 1,\ 4)=24$。底面積 $2\times3=6$、高さ4で、$6\times4=24$ と合います。$\vec c$ がななめにかたむいていても、高さ（$z$ 成分の4）しか体積に関係しないのです。これをスカラー三重積といい、7-4の行列式で、もう一度出会います。
\end{motto}

\begin{motto}[3]
\textbf{電流のまわりの磁場：右ねじの法則}\par
中学校の理科では、まっすぐな導線に電流を流すと、導線のまわりに、輪のような磁場ができることを学びました。その向きは、「電流の向きに右ねじが進むとき、ねじを回す向きが磁場の向き」です（右ねじの法則）。これは④の「回る向きを、軸の向きの矢印で表す」の、ちょうど反対の読み方です。軸の矢印（電流）がわかれば、まわりを回る向き（磁場）がわかります。くわしい式（外積が出てきます）は、物理シリーズの電磁気学と8-10で学びます。
\end{motto}

\begin{column}{台風の渦は、なぜ左回り？}
北半球の台風を上から見ると、雲の渦はいつも左回り（反時計回り）です。地球が回っているので、地面の上を動く空気は、進む向きに対して右へそれていきます。この見かけの力（コリオリの力）は、地球の自転の角速度のベクトル $\vec\omega$ と、空気の速度 $\vec v$ の外積を使って $-2m\,\vec\omega\times\vec v$（$m$ は空気の質量）と表されます。気圧の低い中心へ流れこむ空気が、みな右へそれるので、全体として左回りの渦になります。南半球では、それる向きが反対で、台風（サイクロン）は右回りです。
\end{column}

\begin{senmonsho}[力学・電磁気学の書き方]
\begin{bkhonbun}
点 $\symbf r$ に力 $\symbf F$ がはたらくとき、原点のまわりの力のモーメントは
\[ \symbf N=\symbf r\times\symbf F \]
である。質点の角運動量 $\symbf L=\symbf r\times\symbf p$ について、$\symbf N=\symbf 0$ ならば $\symbf L$ は一定に保たれる。また、磁束密度 $\symbf B$ の中を速度 $\symbf v$ で運動する電荷 $q$ は、次の力を受ける。
\[ \symbf F=q\symbf v\times\symbf B \]
\end{bkhonbun}
\bkyaku{太字はベクトル（4-13）。$\symbf p$ は運動量 $m\symbf v$。「力のモーメントが0なら、角運動量は変わらない」。$\symbf B$ の名前「磁束密度」は、この章の「磁場」のこと（くわしくは電磁気学で）。}
\end{senmonsho}

\begin{column}{わざと「左ねじ」にした、ねじ}
ふつうのねじは右ねじですが、わざと左ねじ（左に回すとしまる）にしてあるねじもあります。自転車の左のペダルは、こぐ回転でゆるまないように、左ねじになっています。プロパンガスのような、燃えるガスの容器の口も、水や空気の道具とまちがえてつながないように、左ねじになっているものがあります。ねじの向きも、安全のための工夫なのです。
\end{column}

\begin{taisetsu}
\def\tsl#1{\makebox[6.4\zw][l]{\sffamily\bfseries\small\color{\bkTT}#1}}
\tsl{モーメント}力×うでの長さ。$M=rF\sin\theta$\\[3pt]
\tsl{平面}$M=xF_y-yF_x$（左回りが＋）\\[3pt]
\tsl{外積の大きさ}$|\vec a\times\vec b|=|\vec a||\vec b|\sin\theta$（平行四辺形の面積）\\[3pt]
\tsl{外積の向き}$\vec a$ にも $\vec b$ にも垂直。$\vec a$ から $\vec b$ へ回す右ねじの進む向き\\[3pt]
\tsl{成分}$\vec a\times\vec b=(a_yb_z-a_zb_y,\ a_zb_x-a_xb_z,\ a_xb_y-a_yb_x)$\\[3pt]
\tsl{性質}$\vec b\times\vec a=-\vec a\times\vec b$。平行なら $\vec 0$\\[3pt]
\tsl{物理}トルク $\vec r\times\vec F$、角運動量 $\vec r\times m\vec v$、磁場からの力 $q\vec v\times\vec B$
\end{taisetsu}

\begin{kotoba}
\small
\bkword{力のモーメント（トルク）}{ものを回す働き。力×うでの長さ。単位 N·m。}
\bkword{右ねじ}{右に回すと進むねじ。回す向きを、軸の向きの矢印で表すときに使う。}
\bkword{外積}{大きさ $|\vec a||\vec b|\sin\theta$、向きは $\vec a$ と $\vec b$ に垂直で右ねじの向きのベクトル。}
\bkword{角運動量}{回る勢いを表すベクトル $\vec r\times m\vec v$。}
\bkword{ローレンツ力}{磁場の中を動く電荷が受ける力 $q\vec v\times\vec B$。}
\end{kotoba}

\bkkigocards{\bkkigocard{$\vec a\times\vec b$}{エー・クロス・ビー}{外積。$\vec a$ と $\vec b$ に垂直なベクトル。大きさ $ab\sin\theta$}{cross product}{4-14}}
{\bkkigocard{$\odot\ \ \otimes$}{手前へ・おくへ}{紙に垂直な向き。$\odot$ は手前へ出る、$\otimes$ はおくへ入る}{out of / into the page}{4-14}}

%% ---------------- 例題 ----------------
\bksection{reidai}{例題でたしかめよう}

\begin{reidai}[シーソーのつり合い]
重さ（重力）\qty{300}{N}の人が、シーソーの支点から左へ\qty{1.2}{m}の所にすわっています。\qty{400}{N}の人は、支点から右へ何mの所にすわればつり合いますか。\par
\kangaekata 左回りの力のモーメント＝右回りの力のモーメント（①）。求める長さを $x$ m とする。\par
\shiki 左回り：$300\times1.2=360$\par
\hspace*{2.6\zw}右回り：$400\times x$\par
\hspace*{2.6\zw}$400x=360$\quad 両辺を400でわって\quad $x=0.9$\par
\kotae \qty{0.9}{m}
\end{reidai}
重い人ほど、支点の近くにすわればよいことがわかります。

\begin{jibun}
重さ\qty{240}{N}の人が支点から左へ\qty{1.5}{m}の所にすわっています。\qty{360}{N}の人は、支点から右へ何mの所にすわればつり合いますか。
\bkjans{$240\times1.5=360$。$360x=360$ より \qty{1.0}{m}}
\end{jibun}

\begin{tsumazuki}{ななめの力で、cos を使ってしまう}
\bkng{うで $r$、力 $F$、なす角 $\theta$ のとき $M=rF\cos\theta$}
\bkok{$M=rF\sin\theta$（うでに直角な部分 $F\sin\theta$ が回す）}
$\theta=0°$（中心に向かって押す）で考えると、$\cos0°=1$ では最大になってしまいます。回らないはずなので、$\sin0°=0$ が正しいとわかります。
\end{tsumazuki}

\begin{reidai}[ななめの力のモーメント]
長さ\qty{0.3}{m}のスパナの先を、スパナとのなす角30°の向きに\qty{50}{N}で引きます。ボルトを回す力のモーメントを求めましょう。\par
\kangaekata $M=rF\sin\theta$（②）。$\theta$ は、うで（スパナ）と力のなす角。\par
\shiki $M=0.3\times50\times\sin30°$\par
\hspace*{2.6\zw}$\sin30°=\dfrac{1}{2}$（4-6）なので\quad $M=0.3\times50\times0.5=7.5$\par
\kotae \bkivoNm{7.5}
\end{reidai}
同じ\qty{50}{N}でも、スパナに直角に引けば $0.3\times50=\bkivoNm{15}$ で、2倍の働きになります。

\begin{jibun}
長さ\qty{0.2}{m}のスパナの先を、スパナとのなす角60°の向きに\qty{60}{N}で引きます。力のモーメントを求めましょう（$\sin60°\approx0.866$）。
\bkjans{$0.2\times60\times0.866\approx\bkivoNm{10.4}$}
\end{jibun}

\begin{reidai}[平面の力のモーメントを成分で]
原点Oを中心に回る板の点P$(0.3,\ 0.4)$（m）に、力 $\vec F=(20,\ 10)$（N）がはたらきます。Oのまわりの力のモーメントを求め、回る向きを答えましょう。\par
\kangaekata $M=xF_y-yF_x$（③）。＋なら左回り、－なら右回り。\par
\shiki $M=0.3\times10-0.4\times20$\par
\hspace*{2.6\zw}$=3-8=-5$\par
\kotae \bkivoNm{5}、右回り（時計回り）
\end{reidai}
上向きの $F_y=10$ は左回りに3、右向きの $F_x=20$ は右回りに8の働きをします。右回りのほうが大きいので、右回りです。

\begin{bkfigure}
\begin{tikzpicture}[x=1mm, y=1mm]
  \draw[bkGray, line width=0.6pt, -{Stealth[length=4pt]}] (-4,0) -- (48,0) node[anchor=west, font=\sffamily\small, text=bkInk] {$x$};
  \draw[bkGray, line width=0.6pt, -{Stealth[length=4pt]}] (0,-4) -- (0,40) node[anchor=south, font=\sffamily\small, text=bkInk] {$y$};
  \draw[bkOrange!70!black, line width=2pt] (0,0) -- (18,24);
  \fill[bkInk] (18,24) circle (1);
  \node[font=\sffamily\footnotesize, anchor=north west, inner sep=1.5pt] at (18.5,23.5) {P$(0.3,\ 0.4)$};
  \draw[bk arrow, bkGreen!70!black, line width=1.6pt] (18,24) -- (18,34);
  \draw[bk arrow, bkBlue, line width=1.6pt] (18,24) -- (38,24);
  \node[font=\sffamily\footnotesize, text=bkGreen!50!black, anchor=east, inner sep=1.5pt] at (17.5,31) {$F_y=10$};
  \node[font=\sffamily\footnotesize, text=bkBlue!80!black, anchor=south, inner sep=1.5pt] at (31,24.4) {$F_x=20$};
  \node[font=\sffamily\footnotesize, anchor=north east, inner sep=1pt] at (-0.5,-0.5) {O};
  \node[font=\sffamily\footnotesize, align=left, anchor=west] at (50,26) {\textcolor{bkGreen!50!black}{$F_y$}：うでの長さ $x=0.3$、左回り\enspace $+0.3\times10=+3$};
  \node[font=\sffamily\footnotesize, align=left, anchor=west] at (50,16) {\textcolor{bkBlue!80!black}{$F_x$}：うでの長さ $y=0.4$、右回り\enspace $-0.4\times20=-8$};
  \node[font=\sffamily\footnotesize, align=left, anchor=west] at (50,6) {合わせて $3-8=-5$（右回り）};
\end{tikzpicture}
\bkfigcaption{例題3の力を、成分ごとに分けて見る。}
\end{bkfigure}

\begin{jibun}
点P$(0.5,\ 0.2)$（m）に、力 $\vec F=(-10,\ 30)$（N）がはたらきます。原点のまわりの力のモーメントを求め、回る向きを答えましょう。
\bkjans{$0.5\times30-0.2\times(-10)=15+2=17$。\bkivoNm{17}、左回り}
\end{jibun}

\begin{reidai}[外積を成分で計算する]
$\vec a=(1,\ 2,\ 0)$、$\vec b=(0,\ 1,\ 3)$ について、\par
（1）$\vec a\times\vec b$ を求め、$\vec a$ とも $\vec b$ とも垂直であることを、内積でたしかめましょう。\par
（2）$\vec b\times\vec a$ を求めましょう。\quad（3）$\vec a\times(2\vec a)$ を求めましょう。\par
\kangaekata 成分の式 $(a_yb_z-a_zb_y,\ a_zb_x-a_xb_z,\ a_xb_y-a_yb_x)$（⑤）。順番を入れかえると符号が反対、平行なら0（⑥）。\par
\shiki（1）$x$ 成分\enspace $2\times3-0\times1=6$\par
\hspace*{4.6\zw}$y$ 成分\enspace $0\times0-1\times3=-3$\par
\hspace*{4.6\zw}$z$ 成分\enspace $1\times1-2\times0=1$\par
\hspace*{4.6\zw}$\vec a$ との内積\enspace $1\times6+2\times(-3)+0\times1=0$\par
\hspace*{4.6\zw}$\vec b$ との内積\enspace $0\times6+1\times(-3)+3\times1=0$\par
\hspace*{2.6\zw}（2）$\vec b\times\vec a=-\vec a\times\vec b=(-6,\ 3,\ -1)$\par
\hspace*{2.6\zw}（3）$2\vec a$ は $\vec a$ と平行なので、$\vec 0$\par
\kotae （1）$(6,\ -3,\ 1)$（内積はどちらも0で、垂直）\quad（2）$(-6,\ 3,\ -1)$\quad（3）$\vec 0=(0,\ 0,\ 0)$
\end{reidai}

\begin{jibun}
$\vec a=(2,\ 1,\ 0)$、$\vec b=(1,\ 0,\ 1)$ について、（1）$\vec a\times\vec b$ を求め、$\vec a$、$\vec b$ との内積が0になることをたしかめましょう。（2）$\vec b\times\vec a$ を求めましょう。（3）$\vec b\times(3\vec b)$ を求めましょう。
\bkjans{（1）\nobreak $(1,\ -2,\ -1)$。内積は $2-2+0=0$、$1+0-1=0$。（2）\nobreak $(-1,\ 2,\ 1)$。（3）\nobreak $\vec 0$}
\end{jibun}

\begin{reidai}[外積で三角形の面積を求める]
3点A$(1,\ 0,\ 0)$、B$(0,\ 2,\ 0)$、C$(0,\ 0,\ 2)$ がつくる三角形の面積を求めましょう。\par
\kangaekata 面積は $\dfrac{1}{2}\bigl|\overrightarrow{\mathrm{AB}}\times\overrightarrow{\mathrm{AC}}\bigr|$（⑥）。\par
\shiki $\overrightarrow{\mathrm{AB}}=(0-1,\ 2-0,\ 0-0)=(-1,\ 2,\ 0)$\par
\hspace*{2.6\zw}$\overrightarrow{\mathrm{AC}}=(0-1,\ 0-0,\ 2-0)=(-1,\ 0,\ 2)$\par
\hspace*{2.6\zw}外積の $x$ 成分\enspace $2\times2-0\times0=4$\par
\hspace*{2.6\zw}$y$ 成分\enspace $0\times(-1)-(-1)\times2=2$\par
\hspace*{2.6\zw}$z$ 成分\enspace $(-1)\times0-2\times(-1)=2$\par
\hspace*{2.6\zw}大きさ\enspace $\sqrt{16+4+4}=\sqrt{24}=2\sqrt6$\par
\hspace*{2.6\zw}面積\enspace $\dfrac{1}{2}\times2\sqrt6=\sqrt6\approx2.45$\par
\kotae $\sqrt6$（約2.45）
\end{reidai}
$\sqrt{24}=\sqrt{4\times6}=2\sqrt6$ は、3-10の√ の計算です。

\begin{jibun}
3点A$(2,\ 0,\ 0)$、B$(0,\ 2,\ 0)$、C$(0,\ 0,\ 2)$ がつくる三角形の面積を求めましょう。
\bkjans{$\overrightarrow{\mathrm{AB}}=(-2,\ 2,\ 0)$、$\overrightarrow{\mathrm{AC}}=(-2,\ 0,\ 2)$。外積 $(4,\ 4,\ 4)$、大きさ $\sqrt{48}=4\sqrt3$。面積 $2\sqrt3\approx3.46$}
\end{jibun}

\begin{reidai}[磁場から導線が受ける力]
$x$ の＋の向きに\qty{3}{A}の電流が流れる長さ\qty{0.2}{m}の導線が、$y$ の＋の向きの\qty{0.5}{T}の磁場の中にあります。導線が受ける力の向きと大きさを求めましょう。\par
\kangaekata $\vec F=I\,\vec L\times\vec B$（⑧）。$\vec L=(0.2,\ 0,\ 0)$、$\vec B=(0,\ 0.5,\ 0)$。\par
\shiki $\vec L\times\vec B=(0\times0-0\times0.5,\ \ 0\times0-0.2\times0,\ \ 0.2\times0.5-0\times0)$\par
\hspace*{4.6\zw}$=(0,\ 0,\ 0.1)$\par
\hspace*{2.6\zw}$\vec F=3\times(0,\ 0,\ 0.1)=(0,\ 0,\ 0.3)$\par
\kotae $z$ の＋の向きに\qty{0.3}{N}
\end{reidai}
$\vec i\times\vec j=\vec k$ から、向きはすぐに $z$ の＋とわかります。フレミングの左手の法則でも同じ向きになることを、自分の左手でたしかめましょう。

\begin{jibun}
$y$ の＋の向きに\qty{2}{A}の電流が流れる長さ\qty{0.5}{m}の導線が、$z$ の＋の向きの\qty{0.4}{T}の磁場の中にあります。導線が受ける力の向きと大きさを求めましょう。
\bkjans{$\vec L\times\vec B=(0.2,\ 0,\ 0)$。$\vec F=(0.4,\ 0,\ 0)$ で、$x$ の＋の向きに\qty{0.4}{N}}
\end{jibun}

\begin{tsumazuki}{左手の法則で、指の役わりを取りちがえる}
\bkng{親指を電流、中指を力にする}
\bkok{中指から「電・磁・力」：中指が電流、人さし指が磁場、親指が力}
取りちがえると、力の向きが反対になります。外積 $I\vec L\times\vec B$ でも計算して、2つの方法で答えをたしかめましょう。
\end{tsumazuki}

\begin{reidai}[トルクのベクトル]
ボルトから見て $\vec r=(0,\ 0.3,\ 0)$（m）の所に、力 $\vec F=(0,\ 0,\ 50)$（N）をかけます。トルク $\vec r\times\vec F$ を求め、ボルトの軸がどの向きかを答えましょう。\par
\kangaekata 外積の成分の式（⑤）で計算する。トルクの矢印の向きが、回す軸の向き（⑦）。\par
\shiki $x$ 成分\enspace $0.3\times50-0\times0=15$\par
\hspace*{2.6\zw}$y$ 成分\enspace $0\times0-0\times50=0$\par
\hspace*{2.6\zw}$z$ 成分\enspace $0\times0-0.3\times0=0$\par
\kotae $(15,\ 0,\ 0)$（\unit[inter-unit-product=\ensuremath{\cdot}]{N.m}）。$x$ 軸の向きの軸のまわりに、右ねじが $x$ の＋へ進む向きに回す
\end{reidai}

\begin{jibun}
ボルトから見て $\vec r=(0.25,\ 0,\ 0)$（m）の所に、力 $\vec F=(0,\ 0,\ -40)$（N）をかけます。トルクを求め、どの向きの軸のまわりに回すかを答えましょう。
\bkjans{$y$ 成分 $0\times0-0.25\times(-40)=10$、ほかは0。$(0,\ 10,\ 0)$。$y$ 軸の向きの軸のまわりに、右ねじが $y$ の＋へ進む向き}
\end{jibun}

\begin{reidai}[いくつもの力がはたらく棒のつり合い]
支点O（原点）で支えた横向きの棒に、図のように力がはたらいています。点 $x=-0.6$（m）に下向き\qty{100}{N}、点 $x=0.4$（m）に下向き\qty{50}{N}。点 $x=1.0$（m）に下向きの力 $F$ を加えて、棒をつり合わせます。$F$ を求めましょう。\par
\noindent\begin{minipage}[c]{0.44\linewidth}
\centering
\begin{tikzpicture}[x=0.54mm, y=1mm]
  \fill[bkGray!50] (0,0) -- (-4.5,-5) -- (4.5,-5) -- cycle;
  \fill[bkOrange!35] (-34,0) rectangle (54,2);
  \draw[bkOrange!70!black, line width=0.5pt] (-34,0) rectangle (54,2);
  \draw[bk arrow, bkBlue, line width=1.4pt] (-30,14) -- (-30,2.5);
  \draw[bk arrow, bkBlue, line width=1.1pt] (20,9) -- (20,2.5);
  \draw[bk arrow, bkRed, line width=1.2pt] (50,11) -- (50,2.5);
  \node[font=\sffamily\footnotesize, anchor=south] at (-30,14) {\qty{100}{N}};
  \node[font=\sffamily\footnotesize, anchor=south] at (20,9) {\qty{50}{N}};
  \node[font=\sffamily\footnotesize, text=bkRed!80!black, anchor=south] at (50,11) {$F$};
  \node[font=\sffamily\footnotesize, anchor=north] at (-30,-1) {$-0.6$};
  \node[font=\sffamily\footnotesize, anchor=north] at (20,-1) {0.4};
  \node[font=\sffamily\footnotesize, anchor=north] at (50,-1) {1.0};
  \node[font=\sffamily\footnotesize, anchor=north] at (0,-5.5) {O};
\end{tikzpicture}
\end{minipage}\hfill
\begin{minipage}[c]{0.53\linewidth}
\kangaekata 力はどれも下向きなので $F_y$ が負、$F_x=0$、$y=0$。1つずつ $M=xF_y-yF_x=xF_y$（③）を計算し、合計を0にする（①のつり合い）。
\end{minipage}\par
\shiki 100 Nの力\enspace $(-0.6)\times(-100)=+60$（左回り）\par
\hspace*{2.6\zw}50 Nの力\enspace $0.4\times(-50)=-20$（右回り）\par
\hspace*{2.6\zw}$F$ の力\enspace $1.0\times(-F)=-F$（右回り）\par
\hspace*{2.6\zw}合計が0なので\enspace $60-20-F=0$、\enspace $F=40$\par
\kotae 下向きに\qty{40}{N}
\end{reidai}
左がわの100 Nは左回り、右がわの2つの力は右回りに回そうとします。左回りの60と、右回りの $20+40=60$ が等しくなって、つり合います。橋やクレーンのつり合いを調べるときも、このように「力のモーメントの合計が0」を使います（物理シリーズ 力-5）。

\begin{jibun}
支点Oで支えた横向きの棒の、点 $x=-0.5$（m）に下向き\qty{120}{N}、点 $x=0.3$（m）に下向き\qty{100}{N}がはたらいています。点 $x=1.5$（m）に下向きの力 $F$ を加えてつり合わせるとき、$F$ を求めましょう。
\bkjans{$(-0.5)\times(-120)=60$、$0.3\times(-100)=-30$、$1.5\times(-F)$。$60-30-1.5F=0$ より $F=\qty{20}{N}$（下向き）}
\end{jibun}

\begin{tsumazuki}{成分の式の $y$ 成分の符号}
\bkng{$y$ 成分を $a_xb_z-a_zb_x$ とする}
\bkok{$y$ 成分は $a_zb_x-a_xb_z$（$z$ と $x$ の順）}
$x$ 成分 $a_yb_z-a_zb_y$ の文字を、$x\to y\to z\to x$ と1つずらして作ります。たすきがけの図で、右下がり－左下がりの順を守りましょう。
\end{tsumazuki}

%% ---------------- 物理・工学 ----------------
\begin{butsuri}
\noindent\begin{minipage}[c]{0.62\linewidth}
\textbf{自転車のペダル}\par\noindent クランク（ペダルのうで）の長さはおよそ\qty{0.17}{m}です。足で真下に\qty{500}{N}でふむと、クランクが水平のときはうでと力が直角で、$0.17\times500=\bkivoNm{85}$ のトルクになります。クランクが真上や真下を向いたときは、うでと力が平行で、トルクは0。ここを「死点」といい、こぐ力が伝わりません。
\end{minipage}\hfill
\begin{minipage}[c]{0.35\linewidth}
\centering
\begin{tikzpicture}[x=1mm, y=1mm]
  \draw[bkGray, line width=0.4pt, dashed] (0,0) circle (14);
  \fill[bkGray!50] (0,0) circle (5);
  \draw[bkInk!70, line width=2.4pt, line cap=round] (0,0) -- (14,0);
  \draw[bkInk!70, line width=2.4pt, line cap=round] (0,0) -- (0,14);
  \fill[bkInk] (0,0) circle (1.2);
  \draw[bk arrow, bkRed, line width=1.4pt] (14,13) -- (14,1.5);
  \draw[bk arrow, bkBlue, line width=1.4pt] (0,27) -- (0,15.5);
  \node[font=\sffamily\footnotesize, text=bkRed!80!black, anchor=north, align=center] at (14,-15.5) {水平：トルク最大};
  \node[font=\sffamily\footnotesize, text=bkBlue!80!black, anchor=west, align=left] at (1.5,24) {真上：\\トルク0（死点）};
\end{tikzpicture}
\end{minipage}
\end{butsuri}

\begin{column}{タイヤのナットは「トルク」でしめる}
自動車のタイヤは、5本ほどのナットで車にとめてあります。ナットは、ゆるすぎると走っているうちに外れ、しめすぎるとボルトがのびて傷みます。そこで、しめる強さは「力」ではなく「トルク」で決められていて、乗用車ではおよそ\bkivoNm{100}です。整備工場では、決めたトルクになると「カチッ」と知らせる、トルクレンチという道具を使います。長さ\qty{0.4}{m}のレンチの先を直角に押すなら、$\bkivoNm{100}\div\qty{0.4}{m}=\qty{250}{N}$ の力（約25\,kgの物の重さ）です。レンチをななめに押すと、同じ力でもトルクは小さくなります（②）。
\end{column}


\begin{butsuri}
\noindent\begin{minipage}[c]{0.62\linewidth}
\textbf{電気モーター}\par\noindent 磁石のあいだに置いた四角いコイルに電流を流すと、向かい合う2本の辺が、ローレンツ力で反対向きに押されます（右の図は、コイルを横から見た断面。$\odot$ と $\otimes$ は電流の向き）。この2つの力が、コイルの軸のまわりにトルクを生み、コイルが回ります。扇風機・電気自動車・ドローンのプロペラも、このしくみで回っています。
\end{minipage}\hfill
\begin{minipage}[c]{0.35\linewidth}
\centering
\begin{tikzpicture}[x=0.8mm, y=0.8mm]
  \fill[bkRed!55] (-26,-8) rectangle (-18,8);
  \node[font=\sffamily\small\bfseries] at (-22,0) {N};
  \fill[bkBlue!45] (18,-8) rectangle (26,8);
  \node[font=\sffamily\small\bfseries] at (22,0) {S};
  \draw[bkGray, line width=0.4pt, dashed] (0,-12) -- (0,12);
  \draw[bkOrange!70!black, line width=1.6pt] (-10,-3) -- (10,-3) -- (10,3) -- (-10,3) -- cycle;
  \bkivoout{(-10,0)}{2mm}
  \bkivoin{(10,0)}{2mm}
  \draw[bk arrow, bkRed, line width=1.3pt] (-10,3.5) -- (-10,12);
  \draw[bk arrow, bkRed, line width=1.3pt] (10,-3.5) -- (10,-12);
  \node[font=\sffamily\footnotesize, text=bkRed!80!black, anchor=south] at (-10,12) {力};
  \node[font=\sffamily\footnotesize, text=bkRed!80!black, anchor=north] at (10,-12) {力};
\end{tikzpicture}
\end{minipage}
\end{butsuri}

\begin{butsuri}
\textbf{スマートフォンの電子コンパス}\par\noindent 半導体の小さな板に電流を流し、そこに磁場がかかると、板の中を流れる電子（p型の半導体では正孔）が $q\vec v\times\vec B$ で横に押されて、板の片がわに集まります。すると、板の両はしに小さな電圧が生まれます（ホール効果）。この電圧をはかって、地球の磁場の向きを知り、地図アプリで北の向きを表示しています。
\end{butsuri}

%% ---------------- 章末問題 ----------------
\bkendproblems
\bkproblemgroup{基本}{3}

\bkproblem 重さ\qty{350}{N}の人が、シーソーの支点から左へ\qty{1.6}{m}の所にすわっています。\qty{400}{N}の人は、支点から右へ何mの所にすわればつり合いますか（\bkreidairef{1}）。
\bkpans{$350\times1.6=560$。$400x=560$ より \qty{1.4}{m}}

\bkproblem 長さ\qty{0.4}{m}のスパナの先を、\qty{30}{N}で引きます。（1）スパナとのなす角30°の向きに引くとき、（2）スパナに直角に引くときの、力のモーメントを求めましょう（\bkreidairef{2}）。
\bkpans{（1）\nobreak $0.4\times30\times0.5=\bkivoNm{6}$。（2）\nobreak $0.4\times30=\bkivoNm{12}$}

\bkproblem 点P$(0.2,\ 0.5)$（m）に、力 $\vec F=(10,\ -20)$（N）がはたらきます。原点のまわりの力のモーメントを求め、回る向きを答えましょう（\bkreidairef{3}）。
\bkpans{$0.2\times(-20)-0.5\times10=-4-5=-9$。\bkivoNm{9}、右回り}

\bkproblem $\vec a=(1,\ 0,\ 2)$、$\vec b=(3,\ 1,\ 0)$ について、$\vec a\times\vec b$ と $\vec b\times\vec a$ を求めましょう。また、$\vec a\times\vec b$ が $\vec a$ に垂直であることを、内積でたしかめましょう（\bkreidairef{4}）。
\bkpans{$\vec a\times\vec b=(-2,\ 6,\ 1)$、$\vec b\times\vec a=(2,\ -6,\ -1)$。内積 $-2+0+2=0$}

\bkproblem 3点A$(1,\ 0,\ 0)$、B$(0,\ 1,\ 0)$、C$(0,\ 0,\ 1)$ がつくる三角形の面積を求めましょう（\bkreidairef{5}）。
\bkpans{$\overrightarrow{\mathrm{AB}}=(-1,\ 1,\ 0)$、$\overrightarrow{\mathrm{AC}}=(-1,\ 0,\ 1)$。外積 $(1,\ 1,\ 1)$、大きさ $\sqrt3$。面積 $\dfrac{\sqrt3}{2}\approx0.866$}

\bkproblemgroup{チャレンジ}{4}

\bkproblem $y$ の＋の向きに\qty{5}{A}の電流が流れる長さ\qty{0.3}{m}の導線が、$x$ の＋の向きの\qty{0.2}{T}の磁場の中にあります。導線が受ける力の向きと大きさを求めましょう（\bkreidairef{6}）。
\bkpans{$\vec L\times\vec B=(0,\ 0,\ -0.06)$。$\vec F=(0,\ 0,\ -0.3)$ で、$z$ の－の向き（下）に\qty{0.3}{N}}

\bkproblem ボルトから見て $\vec r=(0.3,\ 0,\ 0)$（m）の所に、力 $\vec F=(0,\ 20,\ -40)$（N）をかけます。トルク $\vec r\times\vec F$ を求め、その大きさを求めましょう（\bkreidairef{7}）。
\bkpans{$\vec r\times\vec F=(0,\ 12,\ 6)$。大きさ $\sqrt{144+36}=\sqrt{180}\approx\bkivoNm{13.4}$}

\bkproblem 自転車のクランク（長さ\qty{0.17}{m}）が水平から30°下がった所で、真下に\qty{500}{N}でペダルをふみます。クランクと力のなす角は60°です。トルクを求めましょう（$\sin60°\approx0.866$）（\bkreidairef{2}）。
\bkpans{$0.17\times500\times0.866\approx\bkivoNm{73.6}$}

\bkproblem 長さ\qty{0.5}{m}のレンチの先を、レンチとのなす角60°の向きに\qty{200}{N}で押します。トルクを求めましょう。タイヤのナットをしめるのに必要な\bkivoNm{100}にとどきますか（$\sin60°\approx0.866$）（\bkreidairef{2}）。
\bkpans{$0.5\times200\times0.866\approx\bkivoNm{86.6}$。\bkivoNm{100}にとどかない（直角に押せば\bkivoNm{100}になる）}

\begin{furikaeri}
\begin{checklist}
\item 力のモーメントを「力×うでの長さ」で求め、$M=xF_y-yF_x$ で計算できる。\bkfurref{bk:4-14:moment}
\item 外積の向きと大きさがわかり、成分で計算できる。\bkfurref{bk:4-14:gaiseki}
\item トルクや、磁場から受ける力の向きを、外積で求められる。\bkfurref{bk:4-14:torque}
\end{checklist}
\end{furikaeri}

\bkfukidashi{mina}{内積は「同じ向き」、外積は「回す向き」。2つのかけ算で、力の働きが全部見えてきたね。}
